\section{Conclusion} 
PACE builds from individually-calibrated peers in combination with a local protected detector and an expert pool of detectors across diverse anomaly classes, fit and calibrated on benign data only. On average over the 58 target scenarios, use of PACE instead of the mean of seven local detectors gave improvements of between 2.32 and 2.83pp in macro-average performance for each of our four evaluation settings. Improvements were observed at controlled average FPRs $\leq \alpha=0.05$. Ablating expert detectors showed that removing the 'rare' experts had the largest detrimental impact, while some positive detections could be exclusively attributed to local models. Finally, retaining peer-level results separately led to higher detection than either taking the peer max or mean at the cost of a slightly increased average FPR.

Overall, we have demonstrated benefit from both peer and local detection over the benchmark suite. Improvements were largely observed on TON\_IoT and not monotonically tied to the target-fit budget size. The limited target set fit from training-set hold-out, coupled with differing average FPRs across peers, precludes strong conclusions on generalisability. Further validation against unseen data is required before quantifying PACE's utility and false alarm rates.
